% Auto-generated by code/ — DO NOT EDIT BY HAND.
\begin{table}[htbp]
\centering
\caption{How far each prefecture departed from its counterfactual, and whether it came back.}
\label{tab:tohrecovery}
\begin{tabular}{lrrrrcr}
\toprule
 & 2011 & Trough & Final & Mean & Closed & $p$ \\
\midrule
\textit{Real value added} &  &  &  &  &  &  \\
\quad Fukushima & -9.7 & -9.7 (2011) & 0.4 & -1.5 & 2014 & 0.348 \\
\quad Miyagi & -2.6 & -2.6 (2011) & 8.6 & 6.5 & 2012 & 0.043 \\
\quad Iwate & 1.1 & 1.1 (2011) & 9.9 & 8.3 & n/a & 0.087 \\
\midrule
\textit{Mean luminosity per km$^2$} &  &  &  &  &  &  \\
\quad Fukushima & -16.1 & -24.3 (2012) & -18.6 & -20.0 & -- & 0.261 \\
\quad Miyagi & -5.0 & -14.0 (2012) & -13.6 & -11.4 & -- & 0.652 \\
\quad Iwate & -11.1 & -13.1 (2013) & -13.1 & -12.1 & -- & 0.652 \\
\bottomrule
\end{tabular}
\\[2pt]
\parbox[t]{\linewidth}{\footnotesize \emph{Notes:} Every figure is a percentage of the \emph{counterfactual}: what the prefecture would have had. \textit{Trough} is the largest shortfall and the year it occurred; \textit{Final} is the last year of each window (2018 for value added, 2013 for luminosity, which stops before the VIIRS splice); \textit{Mean} averages over the whole post-treatment window. \textit{Closed} is the first year back at or above the counterfactual after falling below it: ``--'' means it had not closed by the end of the window, ``n/a'' that the series never fell below in the first place. $p$ is the in-space placebo probability against all 22 donors; with that many it cannot fall below 0.043.}
\end{table}
